\chapter{The Fast Multipole Method Framework}\label{chapter:fmm-framework}

The \ac{BEM} lends itself naturally to modelling acoustic wave propagation. Transforming the volumetric governing differential equations into boundary integral equations allows finite structures to be treated in an infinite exterior domain without requiring special boundary conditions.
%TODO:  bettter ending for that line please.
However, the required mathematical 
transformation of the problem yields a full, complex, non-Hermitian matrix \cite{chandler-wildeboundaryelementmethods2007}. As the degrees of freedom increase to model larger problem, the memory requirements to hold the large full matrix and the cost to compute the matrix coefficients themselves, scale as $\mathcal{O}(N^2)$, and the computational cost of solving the resulting linear system using direct solvers scales as $\mathcal{O}(N^3)$, where $N$ is the number of degrees of freedom \cite{liufastmultipoleboundary2009}. This makes the \ac{BEM} computationally prohibitive for large-scale problems.

To overcome this computational bottleneck, we must employ a radically different numerical strategy. This chapter details the algorithmic design and the mathematical formulation of the \ac{FMM} framework. The \ac{FMM} decouples far-field interactions through hierarchical spatial decomposition,  truncated analytical basis expansions, and the application of complex mathematical translation operators \cite{coifmanfastmultipolemethod1993}.

This chapter focuses on the mathematical and algorithmic design of the \ac{FMM}  implemented as a part of this thesis. It details the division of the mesh space into a topological tree data structure, the mathematics of the multi-stage multipole translations, and the implementation of near-field interactions.

\section{Fast Multipole Method}\label{sec:fmm}

The \ac{FMM} was introduced in 1987 by Greengard and Rokhlin \cite{greengardfastalgorithmparticle1987} to accelerate large-scale particle simulations involving long-range pairwise potential interactions. Later, it was extended to other problems, including the \ac{BEM} for Acoustic Scattering \cite{rokhlinrapidsolutionintegral1990}. The \ac{FMM} clusters individual boundary element into a hierarchical spatial tree